\noindent{\Large\bfseries Freight network environmental assessment}\par
\noindent{\large Research-design brief: proposed TR-D article}\par
\noindent\small 19 November 2026 --- Working plan, not manuscript findings\par
\medskip

\noindent\textbf{Question and contribution.}
How do parcel-tour construction and emissions formulations affect freight
network environmental estimates in a controlled Dutch application?
The proposed contribution is a reproducible framework and controlled evidence
on those sensitivities. MASS-GT supplies a reference configuration; the SLR
supplies background. Established algorithms are not claimed as new theory.

\smallskip
\noindent\textbf{Three questions.}
(1) Where does NetSy reproduce the reference, and why do residuals remain?
(2) How does Clarke--Wright change parcel operations and emissions under fixed
demand and constraints? (3) How do factor/fleet choices change estimates on
fixed activity and the apparent effect of the scheduling alternative?

\smallskip
\noindent\textbf{Core experiment.}
First distinguish raw source outputs, documented repairs and fixed-trip
assignment controls. Then use four NetSy configurations:
\begin{center}\small
\begin{tabular}{@{}lll@{}}
\toprule Configuration & Parcel scheduling & Emissions configuration \\
\midrule
N00 & Corrected inherited method & Inherited factors \\
N10 & Clarke--Wright savings & Same inherited factors \\
N01 & Same saved activity as N00 & Declared COPERT fleet/factors \\
N11 & Same saved activity as N10 & Same COPERT fleet/factors \\
\bottomrule
\end{tabular}
\end{center}
Hold demand, capacities, service rules, network and supplied speeds fixed.
Reuse identical link activity for factor contrasts. Harmonize fleet and
environmental boundaries where possible; otherwise identify a combined
factor/fleet sensitivity or omit incomparable endpoints. Report scheduling
effects under each factor configuration and their interaction.

\smallskip
\noindent\textbf{Outcomes and randomness.}
Primary: change in network CO\textsubscript{2} over a common modelled period and
carbon boundary, for parcels and the whole freight network. Secondary:
compatible NOx/exhaust PM, spatial patterns, VKT, empty travel, tour counts,
utilization and service feasibility. Pair configurations by saved upstream
realization. A five-replicate pilot will inform a documented precision target
and fixed final replication count; it is not a sufficient-sample claim.
Do not treat links or hours as independent replicates.

\smallskip
\noindent\textbf{Evidence and limits.}
The existing audit contains source comparisons, three upstream NetSy seeds and
one complete seed-level assignment/emissions chain. These are preliminary
records, not new runs at the proposed freeze. The inspected COPERT case uses
assumed fleet technologies, lacks network gradients and has no non-exhaust
input. Grade-informed paths remain an optional extension requiring defensible
data and costs. Source agreement and lower modelled emissions do not establish
measured accuracy, policy effects or geographical transfer.

\smallskip
\noindent\textbf{Next work.}
Freeze versions and inputs; reconcile periods, fleet/factor boundaries and
coverage; refresh controls; lock pilot-informed replication; run the planned
contrasts. Draft Methods and experimental design before Results and Introduction.
Detailed controls, metrics and gates: \texttt{docs/experiment-plan.md}.
The current journal-specific author guide still requires verification.

\smallskip
\noindent\footnotesize
Background: SLR v0.2.0. Proposed model freeze: NetSy v0.5.0
(\texttt{491b1d5}); MASS-GT checkout: \texttt{4f727c3}, with repairs versioned
separately. This is an AI-assisted planning document, not a submitted paper.
